\section{Introduction}\label{sec:intro}

\Cref{fig:triage} shows a customer-support workflow written in OrchLang, the
language this paper is about. It answers a ticket with a small model. When the
customer's account carries a private fraud-risk flag, \code{high\_risk}, it also
asks a larger model to check the reply before the reply is sent. The flag never
reaches a prompt, the workflow's output or the tool that sends the reply. An
information-flow type system with a program-counter label therefore has nothing
to object to: nothing observable at a sink depends on the flag.

The requests do depend on it. For a customer with the flag set, the workflow makes
two requests, one of them to a more expensive model; for everyone else it makes
one. The provider bills per model, so a tenant's invoice shows which customers
were checked. An observer of the encrypted traffic sees request and response
sizes, from which token-length and packet-size attacks have recovered response
text and prompt topics~\cite{weiss2024what,mcdonald2025whisper}, and a
prompt cache shared between users turns the same requests into a timing
channel~\cite{gu2025auditing}. Each of these observers sees a function of the
workflow's requests and responses. If a secret changes the requests, it can change
what they see, even when no secret value crosses any boundary.

\begin{figure}[t]
\lstinputlisting[basicstyle=\ttfamily\footnotesize]{examples/triage.orch}
\vspace{-2pt}
{\scriptsize\ttfamily\raggedright\hangindent=1.5em%% Generated by examples/run.py from `orchc check triage.orch`. Do not edit.
triage.orch:13:3: error [E236] the guard 'high\_risk' depends on a secret and the two arms send different requests, so the trace observer can tell which arm ran; one arm makes 1 top-level request(s), the other 0. then-arm sends [openai/gpt-4o("Check this reply for mistakes: " + reply)] and else-arm sends [nothing]
\par}
\caption{The running example and the compiler's verdict. The secret \code{high\_risk} never
reaches a prompt, an output or a tool, yet it decides whether the second request is sent.
With the header \code{workflow Triage budget 1000000 leaks 1} the same program is accepted
with a warning that it reveals at most one bit (\cref{sec:quantitative}); without the
\code{endorse} on lines 16--17 it is rejected earlier, because text derived from the untrusted
ticket would drive the tool (\cref{sec:language}).}
\label{fig:triage}
\end{figure}

\subsection{Why the classical treatment does not transfer}\label{sec:intro-classical}

This is a resource side channel, and its classical treatment is mature. Ngo,
Dehesa-Azuara, Fredrikson and Hoffmann define resource-aware noninterference and
constant resource, and certify both by combining information-flow typing with
automatic amortised resource analysis (AARA)~\cite{ngo2017verifying}. Çiçek,
Barthe, Gaboardi, Garg and Hoffmann's RelCost bounds the difference in cost
between two executions with relational refinement types~\cite{cicek2017relational}.
Ngo et al.\ quantify over environments that agree on the sizes of secret data, and
RelCost bounds cost differences in terms of sizes and of how much two inputs
differ. Both charge each operation a cost that the program's data determine.

The cost of an LLM call is not determined by sizes. Two requests of exactly the
same length in tokens can receive answers of different lengths, because the model
answers what it is asked. We ran a small instruction-tuned model on fifteen pairs
of chat requests padded to identical token counts, and under a shared sampling
seed all fifteen pairs produced outputs of different lengths
(\cref{sec:eval-models}). Tokenizers disagree with sizes as well: under the
fourteen tokenizers we measured, 84--92\% of random pairs of equal-length strings
receive different token counts. One might hope to repair the classical analysis
with a finer size measure. \Cref{sec:sizes} shows that no size measure suffices:
an analysis that sees prompt text only through a size measure is either unsound
for some provider whose answers depend on content, or it rejects a secret-guarded
branch whose two arms are identical (\cref{thm:size-blind}, mechanised in Coq).
Instantiating Ngo et al.'s framework with a call charged by the sizes of its
request yields such an analysis. RelCost can do better, because it can also state
that two runs with identical inputs have the same cost; charging a call a zero
difference exactly when its two requests are identical is sound, and it is, in
essence, the comparison we develop in \cref{sec:requests}. Ngo et al.'s
quantitative bound fails for a different reason. It counts the possible values of
a cost interval, and an LLM call may return anywhere from zero tokens to its cap,
so the bound is at least $\log_2(\capm+1)$ bits for every workflow that calls a
model, including those that provably leak nothing: 6.0 to 9.4 bits on the eleven
leak-free workflows of our suite, where our bound is zero (\cref{sec:eval-leakage}).

We learned this twice. The first version of our compiler accepted a secret branch
when its arms had equal certified cost bounds, and an external audit showed that
equal maxima do not make equal bills. The repair compared symbolic sizes of
requests; our own audit then showed that equal sizes do not make equal bills
either. \Cref{sec:failed} records both counterexamples, and \cref{thm:size-blind}
explains why every repair of that kind fails.

\subsection{Contributions}\label{sec:contributions}

\begin{itemize}[leftmargin=*]
\item \emph{A size-blindness theorem} (\cref{sec:sizes}). An analysis whose verdict
  depends on the text of a program only through a size measure is unsound for
  content-dependent providers or rejects identical arms, even if soundness is
  demanded only for deterministic providers and only for the total count of output
  tokens. We mechanised the theorem in Coq, and doing so exposed a gap in our own
  pen-and-paper proof.
\item \emph{Content-level request signatures} (\cref{sec:requests}). A branch
  guarded by secret content is resolved for every outcome vector the secret can
  realise, and the resulting request signatures are compared by content. If all
  resolutions agree, the workflow is noninterfering for every provider, tokenizer
  and observer of its transcript (\cref{thm:ni}, mechanised). If $k$ classes remain,
  the channel from the secrets to any such observer has min-capacity at most
  $\log_2 k$, across any number of adaptively chosen invocations
  (\cref{thm:leakage}). Observers that see less than the full transcript admit
  reordering of independent calls, except inside retry bodies (\cref{thm:observers}),
  and a rejection is justified up to dead code (\cref{thm:completeness}).
\item \emph{Token bounds that hold for real tokenizers} (\cref{sec:bounds}). Token
  counts are not subadditive, a response can re-encode to nine tokens per generated
  token, and six of fourteen tokenizers emit more tokens than bytes on some input.
  We therefore compute the guaranteed input bound in bytes and convert it once per
  request through per-tokenizer contracts generated from measurements. The bound
  is proved relative to those contracts (\cref{thm:bound}, mechanised).
\item \emph{An implementation and an evaluation able to fail} (\cref{sec:implementation,sec:evaluation}).
  The compiler \code{orchc} implements both analyses. Its runtime includes a
  provider whose answers depend on the request content, three ways of coupling
  the randomness of two runs, and re-billing with real tokenizers. The evaluation
  includes 30 real workflows ported from 52 candidates under a protocol that fixed
  the labels before any port was written, with a held-out split checked once
  against a frozen compiler.
\end{itemize}

We do not claim that combining information flow with resource analysis is
new~\cite{ngo2017verifying}, nor relational cost analysis~\cite{cicek2017relational},
nor leakage through the token counts, sizes and timing of LLM
services~\cite{weiss2024what,zhang2024time,mcdonald2025whisper,gu2025auditing}.
The flow part of OrchLang is conventional two-point confidentiality and integrity
typing~\cite{denning1976lattice,volpano1996sound}; for agents, CaMeL and FIDES
enforce richer policies at run time~\cite{debenedetti2026defeating,costa2025securing}.
Accepting a secret branch because its arms perform the same observable operations
goes back to Agat's cross-copying transformation~\cite{agat2000transforming}, and
requiring the sequence of observable operations to be independent of the secret is
the program-counter security model of Molnar et al.~\cite{molnar2006program}, with
instructions replaced here by requests. What we add is the proof that, for LLM
calls, ``the same operation'' has to mean the same request content, and an analysis
built on that. Finally, the relational analysis never fires on the 30
real workflows, because none has a secret in its source's threat model; its
practical relevance rests on the threat model, not on deployed workflows
(\cref{sec:eval-real}).

\Cref{sec:setting} fixes the model of providers, randomness and observers.
\Cref{sec:language} introduces OrchLang. \Cref{sec:sizes} presents the two rules
that failed and the theorem behind them, and \cref{sec:requests} the analysis that
replaces them. \Cref{sec:bounds} turns to token bounds, \cref{sec:mechanisation}
to the Coq development and \cref{sec:implementation} to the compiler. We evaluate
in \cref{sec:evaluation} and discuss related work in \cref{sec:related}.
\Cref{sec:conclusion} states the limitations and \cref{sec:conclusion-final}
concludes. The appendices contain the remaining proofs and the corpus of real
workflows.
